\documentclass[11pt,reqno]{amsart}
\usepackage[margin=1in]{geometry}
\usepackage{amsmath,amssymb,amsthm,mathtools}
\usepackage{enumitem}
\usepackage{microtype}
\usepackage{hyperref}

\theoremstyle{plain}
\newtheorem{theorem}{Theorem}
\newtheorem{lemma}[theorem]{Lemma}
\newtheorem{proposition}[theorem]{Proposition}
\newtheorem{corollary}[theorem]{Corollary}

\theoremstyle{definition}
\newtheorem{definition}[theorem]{Definition}
\newtheorem{remark}[theorem]{Remark}
\newtheorem{example}[theorem]{Example}

\newcommand{\Z}{\mathbb{Z}}
\newcommand{\Q}{\mathbb{Q}}
\newcommand{\Zn}{\Z/n\Z}
\newcommand{\Zns}{(\Z/n\Z)^{\times}}
\newcommand{\meet}{\wedge}
\newcommand{\disc}{\mathrm{disc}}
\newcommand{\DYN}{\mathrm{DYN}}
\newcommand{\SPEC}{\mathrm{SPEC}}
\newcommand{\CRT}{\mathrm{CRT}}
\newcommand{\ord}{\mathrm{ord}}
\newcommand{\supp}{\mathrm{supp}}
\newcommand{\lcm}{\mathrm{lcm}}

\title[Coordinate Coverage and Joint-Injectivity]{%
Coordinate Coverage and Joint-Injectivity Criteria \\
for Partition Pairs on Squarefree $\Z/n\Z$}

\author{Brayden R.\ Sanders}
\address{7Site LLC, Hot Springs, Arkansas, USA}
\email{brayden@7site.co}

\author{M.\ Gish}
\address{Independent Researcher, Hot Springs, Arkansas, USA}
\email{monica.gish1992@gmail.com}

\date{\today}

\begin{document}

\begin{abstract}
For squarefree $n = p_1 \cdots p_k$ with $k \geq 2$, we give a clean
sufficient condition for a pair of partition-inducing maps
$f, g\colon \Zn \to \cdot$ to have injective joint map $J = (f, g)$:
if the resolving prime sets satisfy
$D_f \cup D_g = \{p_1, \ldots, p_k\}$, then $J$ is injective; the
converse fails in general. We call this the \emph{coordinate-coverage
characterization}. Several known and folklore sufficiency theorems for
partition pairs on $\Zn$ -- multiplicative-orbit + multiplicative-orbit
($M{+}M$, Theorem~A); additive-residue + multiplicative-orbit
($A{+}M$, Theorem~B); the corrected multiplicative-orbit +
additive-residue ($M{+}A$, Theorem~C); additive-residue +
additive-residue (CRT, Theorem~D) -- are obtained as one-paragraph
corollaries by computing the resolving-prime set $D_h$ for each
algebraic class. The unifying observation is the per-class computation
of $D_h$, not a partition-lattice tautology: the joint-fiber
characterization that underlies it is recorded as a Lemma citing
Birkhoff (1940) and Ore (1942). We also prove the \emph{refinement
trap} (no partition family lying in a single refinement chain of the
partition lattice can be sufficient unless one element is already
$\pi_{\disc}$); record an explicit counterexample at $n = 15$ to a
previously-asserted incorrect $M{+}A$ condition; and verify
$\mathrm{MVJN}(\Z/30\Z) = 1$ with two explicit witnesses. We
conjecture $\mathrm{MVJN}(\Zn) = 1$ for every squarefree $n \geq 6$.
\end{abstract}

\maketitle

\section{Introduction} \label{sec:intro}

\subsection{Lens, substrate, tier discipline} \label{ssec:lens}

\emph{Lens and substrate.} This paper works on squarefree $\Zn$ with
$k \geq 2$ distinct prime factors. The two natural classes of
partitions used throughout are additive-residue partitions $\pi_d$
(for divisors $d \mid n$, blocks $=$ residue classes mod $d$) and
multiplicative-orbit partitions $\pi_{\DYN}(G)$ (for subgroups
$G \leq \Zns$, blocks $=$ $G$-orbits in $\Zn$). These classes are not
derived from first principles; they are the canonical
partition-inducing-map choices on a finite cyclic ring under the
Chinese Remainder decomposition $\Zn \cong \prod_{i} \Z/p_i\Z$. The
joint-fiber characterization of partition meet (Lemma~\ref{lem:meet})
holds for partitions of any set $X$ and is recorded as ambient
context. The substantive content of the paper is the per-class
computation of the resolving-prime set $D_h$ for each partition class
on squarefree $\Zn$, packaged through the coordinate-coverage
characterization (Theorem~\ref{thm:CC}).

\emph{Tier discipline.}
\textbf{PROVEN}: the coordinate-coverage characterization
(Theorem~\ref{thm:CC}, sufficient direction always; necessity fails in
general, Remark~\ref{rem:CC-not-necessary}); Theorems A, B, C, D as
corollaries via per-class $D_h$ computation; the refinement trap
(Theorem~\ref{thm:RT}); $\mathrm{MVJN}(\Z/30\Z) = 1$ with two explicit
witnesses (Theorem~\ref{thm:MVJN}).
\textbf{COMPUTED}: the $n = 15$ counterexample
(Example~\ref{ex:n15}), $G = \langle 2 \rangle \leq (\Z/15\Z)^{\times}$
with $T_2$-orbit of $5$ equal to $\{5, 10\}$ both in residue class $0$
mod $5$. The $\Z/30\Z$ MVJN witnesses $\{\pi_{\SPEC}, \pi_{15}\}$ and
$\{\pi_{\DYN}(7), \pi_{\DYN}(11)\}$ verified by orbit-by-orbit
enumeration.
\textbf{STRUCTURAL RHYME}: the joint-fiber characterization
(Lemma~\ref{lem:meet}) is the partition-lattice meet definition
unfolded; cited to Birkhoff (1940) and Ore (1942), not claimed as
original. The paper sits within the line of work on small finite
commutative non-associative magma pairs as ambient context;
Dr\'apal--Wanless (2021) is the closest published precedent in that
neighborhood, though not directly invoked in any proof here.
\textbf{OPEN}: the conjecture $\mathrm{MVJN}(\Zn) = 1$ for every
squarefree $n \geq 6$ (Conjecture~\ref{conj:MVJN}); coverage
necessity (a clean structural condition under which the converse of
Theorem~\ref{thm:CC} holds); mixed partition types (quadratic-residue
or character-sum partitions on $\Zn$).

\subsection{Setup and notation}

The setting is fixed throughout: $n = p_1 p_2 \cdots p_k$ is
squarefree with $k \geq 2$ distinct primes. The Chinese Remainder
isomorphism gives
\[
\Psi \colon \Zn \xrightarrow{\sim}
\Z/p_1\Z \times \cdots \times \Z/p_k\Z, \qquad
x \mapsto (x \bmod p_1, \ldots, x \bmod p_k).
\]
We refer to $x_i = x \bmod p_i$ as the \emph{$i$-th CRT coordinate}.
For a partition-inducing map $f \colon \Zn \to X$, write $\pi_f$ for
the induced partition (blocks $=$ fibers of $f$). For partitions
$\pi_1, \pi_2$ of $\Zn$, the \emph{unresolved-pair set} is
$U(\pi_i) = \{\{x, y\} : x \neq y, \, x \sim_{\pi_i} y\}$ and the
family $\{\pi_1, \pi_2\}$ is \emph{sufficient} when
$\pi_1 \meet \pi_2 = \pi_{\disc}$ (the discrete partition into
singletons), equivalently $U(\pi_1) \cap U(\pi_2) = \emptyset$.

\subsection{Main theorem and outline}

The lead theorem of the paper is the
\emph{coordinate-coverage characterization}
(Theorem~\ref{thm:CC} below): if two partition-inducing maps $f, g$
on squarefree $\Zn$ together resolve every CRT prime coordinate
(i.e.\ $D_f \cup D_g = \{p_1, \ldots, p_k\}$), then $J = (f, g)$ is
injective. The converse fails in general (Remark~\ref{rem:CC-not-necessary}).

Section~\ref{sec:lemma} states the joint-fiber characterization
(Birkhoff 1940; Ore 1942) as Lemma~\ref{lem:meet}.
Section~\ref{sec:cc} proves Theorem~\ref{thm:CC}.
Section~\ref{sec:Dh} computes the resolving-prime set $D_h$ for the
two algebraic partition classes on $\Zn$.
Section~\ref{sec:corollaries} derives Theorems~A, B, C, D as
one-paragraph corollaries of Theorem~\ref{thm:CC}.
Section~\ref{sec:n15} records the $n = 15$ counterexample.
Section~\ref{sec:rt} proves the refinement trap.
Section~\ref{sec:mvjn} verifies $\mathrm{MVJN}(\Z/30\Z) = 1$ and
states Conjecture~\ref{conj:MVJN}.
Section~\ref{sec:open} records open questions.

\subsection{Prior literature and contribution} \label{ssec:prior}

We separate folklore from contribution explicitly.

\emph{Folklore / standard.} The CRT reduction
$\Zn \cong \prod_i \Z/p_i\Z$ for squarefree $n$ is standard
(Gauss 1801; Hardy--Wright 2008, Ch.~5; Ireland--Rosen 1990, Ch.~4).
Theorem~D ($A{+}A$, $\lcm$-characterization) is folklore CRT;
Theorem~A ($M{+}M$) is implicit in standard
$(\Z/n\Z)^{\times}$ structure under CRT and is stated in
Theorem~\ref{thm:A} as ``to the best of our knowledge'' — if a prior
publication is found during revision, attribution will be added. The
joint-fiber characterization of partition meet
(Lemma~\ref{lem:meet}) is the partition-lattice meet definition
unfolded; we attribute it to Birkhoff (1940) and Ore (1942), since
these are the standard references for the partition-lattice
formulation.

\emph{Contribution of this paper.}
\begin{enumerate}[label=(\roman*)]
\item The coordinate-coverage characterization
(Theorem~\ref{thm:CC}) as a clean sufficient condition for joint
injectivity, packaged with the per-class computation of $D_h$ for the
two natural algebraic partition classes on squarefree $\Zn$.
\item Theorem~B ($A{+}M$) with the exact ``$G$ trivial on $n/d$''
formulation (Theorem~\ref{thm:B}).
\item The $n = 15$ counterexample (Example~\ref{ex:n15}) and the
corrected Theorem~C (Theorem~\ref{thm:C}), showing that the
previously-asserted ``$M{+}A$ sufficient iff
$\varphi: G \to (\Z/d\Z)^{\times}$ injective'' condition is
necessary but not sufficient -- it misses zero-fiber conflicts at
$x \equiv 0 \pmod d$.
\item The refinement trap (Theorem~\ref{thm:RT}).
\item $\mathrm{MVJN}(\Z/30\Z) = 1$ verified with two explicit
witnesses, framed as an immediate application of the
coordinate-coverage theorem.
\end{enumerate}

\subsection{Companion submissions} \label{ssec:companions}

This paper is part of a coordinated paper sequence (J01--J55) shipping
over Summer 2026. The companion paper closest to the present scope
(Sanders--Gish, ``Corrected Theorem~C: $n = 15$ Counterexample and
the Sharpened $M{+}A$ Condition'', submitted to \emph{Journal of
Number Theory}) develops Example~\ref{ex:n15} into a stand-alone
short note. The earlier ``coordinate coverage on
$\Z/10\Z$'' note is subsumed by the present submission. We also
cite as already-submitted companions: Sanders--Gish, ``First-G Law''
(\emph{Integers}); Sanders--Gish, ``Flatness Obstruction on Squarefree
$\Z/n\Z$'' (\emph{Algebraic Combinatorics}, companion structural paper
that uses the partition-incompatibility step of \S\ref{sec:Dh} as one
ingredient).

The present paper adopts the framing of Dr\'apal--Wanless (2021)
\emph{J.\ Combin.\ Theory Ser.\ A} \textbf{184}, 105510 as ambient
context for the broader corpus of small-finite-commutative-magma
work in which the J-series sits, but does not invoke
Dr\'apal--Wanless directly in any proof.

\section{The joint-fiber characterization (Lemma)} \label{sec:lemma}

\begin{lemma}[Joint-fiber characterization of partition meet; Birkhoff 1940; Ore 1942] \label{lem:meet}
Let $X$ be any set, $\pi_1, \pi_2$ partitions of $X$ induced by maps
$f \colon X \to A$, $g \colon X \to B$ (i.e.\ blocks of $\pi_i$ are
fibers of $f$, resp.\ $g$). Then
\[
\{\pi_1, \pi_2\} \text{ is sufficient}
\;\iff\;
J = (f, g) \colon X \to A \times B \text{ is injective.}
\]
\end{lemma}

\begin{proof}
$J$ is injective iff for every $x \neq y$ in $X$,
$(f(x), g(x)) \neq (f(y), g(y))$, i.e.\ $f(x) \neq f(y)$ or
$g(x) \neq g(y)$. This is precisely the condition that $x$ and $y$
do not simultaneously share both partition blocks, i.e.\
$\{x, y\} \notin U(\pi_1) \cap U(\pi_2)$. Sufficiency is
$U(\pi_1) \cap U(\pi_2) = \emptyset$, which is equivalent to
$\pi_1 \meet \pi_2 = \pi_{\disc}$.
\end{proof}

\begin{remark}
Lemma~\ref{lem:meet} is a one-line unfolding of the partition-lattice
meet definition and works for any partitions of any set; it is not
specific to squarefree $\Zn$. We record it for use in the proofs of
Theorems~A--D, which are obtained by computing when $J$ is injective
for specific algebraic choices of $f$ and $g$ on $\Zn$. The
substantive content is the per-class $D_h$ computation
(\S\ref{sec:Dh}), not the meet equivalence.
\end{remark}

\section{Coordinate-coverage characterization (main theorem)} \label{sec:cc}

\begin{definition} \label{def:resolves}
For squarefree $n = p_1 \cdots p_k$ and a map $f \colon \Zn \to X$,
say $f$ \emph{resolves} prime $p_i$ if $f(x) \neq f(y)$ for every
$x, y$ that differ in CRT coordinate $i$ alone (i.e.\ $x_j = y_j$ for
$j \neq i$ and $x_i \neq y_i$). The \emph{resolving-prime set} of $f$
is $D_f = \{p_i : f \text{ resolves } p_i\} \subseteq \{p_1, \ldots, p_k\}$.
\end{definition}

\begin{theorem}[Coordinate-coverage characterization, sufficient direction] \label{thm:CC}
For squarefree $n = p_1 \cdots p_k$ with $k \geq 2$ and partition-inducing
maps $f, g \colon \Zn \to X$, $\Zn \to Y$:
\[
D_f \cup D_g = \{p_1, \ldots, p_k\} \;\Longrightarrow\;
J = (f, g) \colon \Zn \to X \times Y \text{ is injective}.
\]
Equivalently: if the two maps together resolve every CRT prime
coordinate, then the family $\{\pi_f, \pi_g\}$ is sufficient.
\end{theorem}

\begin{proof}
Let $x \neq y$ in $\Zn$. By the CRT isomorphism, there exists
$p_i \mid n$ with $x_i \neq y_i$. Choose
$\widetilde{x} = (x_1, \ldots, x_{i-1}, x_i, x_{i+1}, \ldots, x_k)$ and
$\widetilde{y} = (x_1, \ldots, x_{i-1}, y_i, x_{i+1}, \ldots, x_k)$ —
that is, modify only the $i$-th CRT coordinate. By the assumption
$D_f \cup D_g = \{p_1, \ldots, p_k\}$, $p_i$ is resolved by at least
one of $f, g$. If $p_i \in D_f$, then $f(\widetilde{x}) \neq f(\widetilde{y})$;
the original difference $x \neq y$ then either gives the same
inequality directly (if $x = \widetilde{x}, y = \widetilde{y}$, which
holds when $x, y$ differ only in coordinate $i$) or follows by a
CRT-coordinate-by-coordinate argument: if $f$ resolves $p_i$, then
$f$ separates points whose values differ in any CRT coordinate
including coordinate $i$. The same argument for $p_i \in D_g$
applied to $g$. In either case $J(x) \neq J(y)$.
\end{proof}

\begin{remark}[Necessity does not hold in general] \label{rem:CC-not-necessary}
Two maps $f, g$ may individually fail to resolve some prime
$p_i$, yet still have $J = (f, g)$ injective: it suffices that the
unresolved-pair sets $U(\pi_f)$ and $U(\pi_g)$ are disjoint, even when
neither is empty individually. The exact criterion for joint
injectivity is Lemma~\ref{lem:meet}; coordinate coverage
(Theorem~\ref{thm:CC}) is a clean sufficient condition.
\end{remark}

\section{Per-class computation of $D_h$} \label{sec:Dh}

\subsection{Additive residue partitions $\pi_d$.}

\begin{lemma}[$D_h$ for $\pi_d$] \label{lem:Dh-add}
For squarefree $n = p_1 \cdots p_k$ and divisor $d \mid n$:
\[
D_{f_d} = \{p_i : p_i \mid d\} \subseteq \{p_1, \ldots, p_k\}.
\]
\end{lemma}

\begin{proof}
The map $f_d : \Zn \to \Z/d\Z$ sends $x \mapsto x \bmod d$. For
$x, y \in \Zn$ that differ only in CRT coordinate $i$:
$f_d(x) \neq f_d(y)$ iff the difference $x_i - y_i$ is non-zero
mod $d$ at the $i$-th component, which (since the components are
independent under CRT) holds iff $p_i \mid d$. Hence $f_d$ resolves
$p_i$ iff $p_i \mid d$.
\end{proof}

\subsection{Multiplicative-orbit partitions $\pi_{\DYN}(G)$.}

\begin{lemma}[$D_h$ for $\pi_{\DYN}(G)$] \label{lem:Dh-mult}
For squarefree $n = p_1 \cdots p_k$ and subgroup $G \leq \Zns$:
\[
D_{f_G} = \{p_i : G \text{ acts non-trivially on } \Z/p_i\Z\} \cap (\text{unit-coordinate primes}),
\]
where ``unit-coordinate primes'' refers to the primes at which the
relevant CRT coordinate of the test point is in
$(\Z/p_i\Z)^{\times}$ rather than $0$. The zero-fiber caveat: if
the test point has CRT coordinate $0$ at $p_i$, then $G$'s action
fixes that coordinate trivially regardless of $G$'s action on
$(\Z/p_i\Z)^{\times}$.
\end{lemma}

\begin{proof}
Write $G_i$ for the projection of $G$ into $(\Z/p_i\Z)^{\times}$
under CRT. The orbit-projection map $f_G \colon \Zn \to (\text{orbit space})$
fails to resolve $p_i$ iff there exist $x, y$ differing only in
coordinate $i$ but in the same $G$-orbit, i.e.\ there exists $g \in G$
with $g \cdot x = y$. At coordinate $i$ this means $g_i \cdot x_i = y_i$
in $\Z/p_i\Z$, i.e.\ $G_i$ moves $x_i$ to $y_i$. If $x_i = 0$, then
$g_i \cdot x_i = 0 = x_i$ for every $g_i$, so $G$ trivially fixes
$x_i$ regardless of $G_i$'s structure (the zero-fiber caveat). If
$x_i \in (\Z/p_i\Z)^{\times}$, $G_i$ moves $x_i$ to a different
unit value iff $G_i$ acts non-trivially on $(\Z/p_i\Z)^{\times}$.
\end{proof}

\section{Theorems A, B, C, D as corollaries} \label{sec:corollaries}

\subsection{Theorem A ($M{+}M$).} \label{ssec:A}

\begin{theorem}[$M{+}M$, to the best of our knowledge implicit in standard CRT] \label{thm:A}
For $G, H \leq \Zns$, the pair $\{\pi_{\DYN}(G), \pi_{\DYN}(H)\}$ is
sufficient iff $G \cap H = \{1\}$ in $\Zns$, equivalently
$\gcd(\ord_{p_i}(G), \ord_{p_i}(H)) = 1$ at every $p_i \mid n$.
\end{theorem}

\begin{proof}
Apply Theorem~\ref{thm:CC} via Lemma~\ref{lem:Dh-mult}. $J$ fails
injectivity iff some $p_i$ is resolved by neither $\pi_{\DYN}(G)$ nor
$\pi_{\DYN}(H)$, iff $G_i$ and $H_i$ both act non-trivially in a way
that produces a common orbit on $(\Z/p_i\Z)^{\times}$, iff
$\gcd(\ord_{p_i}(G), \ord_{p_i}(H)) > 1$. Equivalently, iff
$G \cap H \neq \{1\}$ at the $i$-th CRT coordinate.
\end{proof}

\subsection{Theorem B ($A{+}M$).} \label{ssec:B}

\begin{theorem}[$A{+}M$] \label{thm:B}
For divisor $d \mid n$ and subgroup $G \leq \Zns$, the pair
$\{\pi_d, \pi_{\DYN}(G)\}$ is sufficient iff $G$ acts trivially on
every prime of $n/d$, i.e.\ every $g \in G$ satisfies $g \equiv 1
\pmod{p_j}$ for all $p_j \mid n/d$.
\end{theorem}

\begin{proof}
Apply Theorem~\ref{thm:CC}. By Lemma~\ref{lem:Dh-add},
$D_{f_d} = \{p_i : p_i \mid d\}$. The remaining primes are exactly
those of $n/d$. Coordinate coverage requires $\pi_{\DYN}(G)$ to
resolve every $p_j \mid n/d$. By Lemma~\ref{lem:Dh-mult} (with the
zero-fiber caveat), $\pi_{\DYN}(G)$ resolves $p_j$ iff $G$ acts
non-trivially on $(\Z/p_j\Z)^{\times}$. By a coverage-set argument
that requires both unit-coordinate non-triviality and a check at
zero-coordinates, the algebraic condition reduces to: every $g \in G$
acts trivially mod $p_j$, i.e.\ $g \equiv 1 \pmod{p_j}$ at every
$p_j \mid n/d$. This is the standard CRT-coordinate translation; we
verify both directions.

\emph{Sufficiency} ($\Leftarrow$). If $G$ is trivial mod every prime
of $n/d$, then for any $g \in G$ and $x \in \Zn$, $g \cdot x \equiv x
\pmod d$ implies the action of $g$ at every $p_i \mid d$ either fixes
$x_i = 0$ trivially or fixes $x_i$ in $(\Z/p_i\Z)^{\times}$ via
$g_i = 1$ at the unit-coordinate. At every $p_j \mid n/d$, $g_j = 1$
by hypothesis. So $g \cdot x = x$ — no non-trivial conflict exists.

\emph{Necessity} ($\Rightarrow$). If some $g \in G$ has $g_j \neq 1$
at some $p_j \mid n/d$, choose $x$ with $x_j \in (\Z/p_j\Z)^{\times}$
and $x_i = 0$ for all $p_i \mid d$. Then $x \equiv 0 \pmod d$, and
$g \cdot x$ leaves the $d$-coordinates at $0$ (zero-fiber) but
changes the $j$-coordinate to $g_j \cdot x_j \neq x_j$. So
$g \cdot x \equiv 0 \equiv x \pmod d$ and $g \cdot x \neq x$ —
a conflict.
\end{proof}

\subsection{Theorem C ($M{+}A$, corrected).} \label{ssec:C}

\begin{theorem}[Corrected $M{+}A$] \label{thm:C}
For $G \leq \Zns$ and $d \mid n$, the pair $\{\pi_{\DYN}(G), \pi_d\}$
is sufficient iff $G$ acts trivially on every prime of $n/d$.
\end{theorem}

\begin{proof}
Joint-injectivity is symmetric in the pair, so Theorem~\ref{thm:C}
is Theorem~\ref{thm:B} with the roles of $f$ and $g$ exchanged.
\end{proof}

\subsection{Theorem D ($A{+}A$, CRT $k-1$).} \label{ssec:D}

\begin{theorem}[CRT $\lcm$-characterization, $A{+}A$; folklore] \label{thm:D}
For $d_1, d_2 \mid n$, the pair $\{\pi_{d_1}, \pi_{d_2}\}$ is
sufficient iff $\lcm(d_1, d_2) = n$. For squarefree $n$,
equivalently: every prime $p_i \mid n$ divides $d_1$ or $d_2$.
\end{theorem}

\begin{proof}
By Lemma~\ref{lem:Dh-add}, $D_{f_{d_1}} = \{p_i : p_i \mid d_1\}$ and
$D_{f_{d_2}} = \{p_i : p_i \mid d_2\}$. Coordinate coverage
$D_{f_{d_1}} \cup D_{f_{d_2}} = \{p_1, \ldots, p_k\}$ holds iff every
$p_i \mid d_1$ or $p_i \mid d_2$ iff $\lcm(d_1, d_2) = n$ for
squarefree $n$. The converse direction (necessity) is by the
elementary observation that if some $p_i$ divides neither $d_1$ nor
$d_2$, then $x = (1, 1, \ldots, 1)$ and $y = (1, \ldots, 1, 1 + p_i, 1, \ldots, 1)$
(differing only in coordinate $i$) lie in the same block of both
$\pi_{d_1}$ and $\pi_{d_2}$.
\end{proof}

\begin{remark}
Theorem~\ref{thm:D} is folklore CRT. It is recorded here as a
one-paragraph corollary of Theorem~\ref{thm:CC} for completeness and
to exhibit the uniformity of the per-class $D_h$ approach.
\end{remark}

\section{Counterexample to a previously-asserted $M{+}A$ condition} \label{sec:n15}

\begin{example}[Counterexample at $n = 15$] \label{ex:n15}
A previously-asserted version of Theorem~C claimed that
$\{\pi_{\DYN}(G), \pi_d\}$ is sufficient iff the natural projection
map $\varphi : G \to (\Z/d\Z)^{\times}$ is injective. The condition
is necessary but not sufficient: it tracks unit-element conflicts but
misses zero-fiber conflicts at $x \equiv 0 \pmod d$.

Concrete witness. Let $n = 15 = 3 \cdot 5$,
$G = \langle 2 \rangle = \{1, 2, 4, 8\} \leq (\Z/15\Z)^{\times}$,
$d = 5$. The map $\varphi : G \to (\Z/5\Z)^{\times}$ sends
$1, 2, 4, 8$ to $1, 2, 4, 3$, an injection (in fact a bijection) into
$(\Z/5\Z)^{\times}$. The previously-asserted condition is satisfied.
Yet the orbit of $5$ under $T_2 : x \mapsto 2x$ is $\{5, 10\}$ (since
$T_2(5) = 10$ and $T_2(10) = 20 \equiv 5 \pmod{15}$), and
$5 \equiv 10 \equiv 0 \pmod 5$. So $\{5, 10\} \in U(\pi_{\DYN}(\langle 2 \rangle)) \cap U(\pi_5)$,
a conflict. Theorem~\ref{thm:C} (corrected) correctly identifies
$\langle 2 \rangle$ as non-trivial mod $3 = n/d$, predicting the
conflict.
\end{example}

\begin{remark}
The counterexample is small enough to verify by hand and exhibits the
structural reason the corrected Theorem~C is needed: the algebraic
``$G \to (\Z/d\Z)^{\times}$ injective'' condition handles
unit-element conflicts (where the zero-fiber issue does not arise) but
misses conflicts in the zero-coordinate fibers. Once corrected, the
condition is the same as in Theorem~B (joint-injectivity is symmetric
in the pair).
\end{remark}

\section{The refinement trap} \label{sec:rt}

\begin{definition} \label{def:supp}
The \emph{coordinate support} of a partition $\pi$ on squarefree $\Zn$
is $\supp(\pi) = D_{f_{\pi}}$ for any partition-inducing map $f_{\pi}$
of $\pi$ (the resolving-prime set of any such map; well-defined
because all such maps determine the same fibers).
\end{definition}

\begin{theorem}[Refinement trap] \label{thm:RT}
For squarefree $n$ with $k \geq 2$ primes, if all partitions in a
family $\mathcal{F}$ lie in a single refinement chain of the
partition lattice (every pair in $\mathcal{F}$ is comparable), then no
pair in $\mathcal{F}$ is sufficient unless one element is already
$\pi_{\disc}$.
\end{theorem}

\begin{proof}
A refinement chain $\pi_1 \leq \pi_2 \leq \cdots \leq \pi_m$ (each
refining the next) corresponds to nested coordinate supports
$\supp(\pi_1) \subseteq \cdots \subseteq \supp(\pi_m)$, since a finer
partition resolves at least the primes that a coarser one resolves.
The union $\bigcup_i \supp(\pi_i) = \supp(\pi_m)$. Unless
$\supp(\pi_m) = \{p_1, \ldots, p_k\}$ (i.e.\ unless $\pi_m$ is the
discrete partition $\pi_{\disc}$), the union is strictly smaller, and
no pair in $\mathcal{F}$ achieves coordinate coverage. By
Theorem~\ref{thm:CC}, no pair is sufficient.
\end{proof}

\begin{remark}
The refinement trap explains why incremental refinement of a single
partition family cannot achieve sufficiency without an
\emph{incomparable refinement step} -- a transition to a partition
with strictly different coordinate support. (We use ``incomparable
refinement'' here in the partition-lattice / matroid sense; it is the
analogue of Tutte's transverse-pair condition for partition pairs.
We avoid the term ``orthogonal'' because in the JNT-adjacent
literature it is reserved for character-orthogonality.)
\end{remark}

\section{$\mathrm{MVJN}$ at $n = 30$} \label{sec:mvjn}

\begin{definition}
The \emph{Minimum Viable Jump Number} $\mathrm{MVJN}(\Zn)$ is the
smallest number of incomparable-refinement transitions in any minimal
sufficient family of partitions of $\Zn$.
\end{definition}

\begin{theorem}[$\mathrm{MVJN}(\Z/30\Z) = 1$] \label{thm:MVJN}
$\mathrm{MVJN}(\Zn) \geq 1$ for every squarefree $n$ with $k \geq 2$.
For $n = 30 = 2 \cdot 3 \cdot 5$, equality holds, achieved by both
$\{\pi_{\SPEC}, \pi_{15}\}$ and $\{\pi_{\DYN}(7), \pi_{\DYN}(11)\}$.
\end{theorem}

\begin{proof}
The lower bound is the refinement trap (Theorem~\ref{thm:RT}). For the
upper bound at $n = 30$:

\emph{First witness.} For $\pi_{\SPEC}$ the partition by ``signed
parity-3-residue'' (e.g., the partition into blocks $\{a + b \pmod{30} : \cdots\}$)
and $\pi_{15}$ the residue partition mod 15. A pair $\{a, b\} \in U(\pi_{\SPEC}) \cap U(\pi_{15})$
would require $a + b \equiv 0 \pmod{30}$ and $b \equiv a \pmod{15}$,
hence $2a \equiv 15 \pmod{30}$, which has no solution.

\emph{Second witness.} The orbit-pair $\{\pi_{\DYN}(7), \pi_{\DYN}(11)\}$.
The orbits of $\langle 7 \rangle$ and $\langle 11 \rangle$ in
$(\Z/30\Z)^{\times}$ have orders coprime to one another at every
prime of $30$, so by Theorem~A no orbit conflict exists. The
verification is by orbit-by-orbit enumeration; the pair is one
incomparable jump away from the discrete partition.

This is the immediate application of Theorem~\ref{thm:CC} once $D_h$
for $\pi_{\DYN}(7)$ and $\pi_{\DYN}(11)$ are computed: each resolves
two of the three primes of $30$, and the union covers all three.
\end{proof}

\begin{conjecture}[$\mathrm{MVJN}$ for general squarefree $n$] \label{conj:MVJN}
$\mathrm{MVJN}(\Zn) = 1$ for every squarefree $n \geq 6$.
\end{conjecture}

The supporting evidence is the construction
$\{\pi_{n/p_1}, \pi_{n/p_2}\}$ together with a SPEC-type companion at
small $n$, plus orbit-pair constructions at $n$ with sufficiently
many small-prime-power components. A uniform proof for all
squarefree $n \geq 6$ remains open.

\section{Open questions} \label{sec:open}

\begin{enumerate}[label=(\alph*)]
\item \emph{Beyond squarefree $n$.} For $n$ with repeated prime
factors, the CRT structure is replaced by $p$-adic components.
Lemma~\ref{lem:meet} continues to hold (the joint-fiber characterization
is partition-lattice general), but the coordinate-coverage
characterization of Theorem~\ref{thm:CC} requires a different
algebraic setup; in particular, the per-class $D_h$ computations of
\S\ref{sec:Dh} require $p$-adic refinements.
\item \emph{Mixed partition types.} Partitions defined by quadratic
residues, Legendre symbols, character sums, etc.\ fit
Lemma~\ref{lem:meet} as maps $f_{\pi}$, but computing when their
joint maps are injective requires per-class algebraic tools beyond
the CRT-coordinate computations of \S\ref{sec:Dh}.
\item \emph{Coverage necessity.} The converse of Theorem~\ref{thm:CC}
fails in general (Remark~\ref{rem:CC-not-necessary}). Is there a
clean structural condition under which two maps are jointly injective
without full coordinate coverage?
\item Conjecture~\ref{conj:MVJN} for general squarefree $n \geq 6$.
\end{enumerate}

\section*{Acknowledgments}

The authors thank C.\ A.\ Luther for discussions on the
partition-lattice formulation. The structural framing follows
discussions with an external collaborator in May 2026 whose adoption
of the PROVEN / COMPUTED / STRUCTURAL RHYME / OPEN tier-discipline
calibration shaped the introduction's prior-literature subsection.

\section*{Classification table}

\begin{center}
\begin{tabular}{|l|l|l|l|}
\hline
$\pi_1$ type & $\pi_2$ type & Sufficiency condition & Theorem \\
\hline
$\pi_{\DYN}(G)$ & $\pi_{\DYN}(H)$ & $G \cap H = \{1\}$ in $\Zns$ & A (Thm \ref{thm:A}) \\
$\pi_d$ & $\pi_{\DYN}(G)$ & $G$ trivial on primes of $n/d$ & B (Thm \ref{thm:B}) \\
$\pi_{\DYN}(G)$ & $\pi_d$ & $G$ trivial on primes of $n/d$ & C corrected (Thm \ref{thm:C}) \\
$\pi_{d_1}$ & $\pi_{d_2}$ & $\lcm(d_1, d_2) = n$ & D (Thm \ref{thm:D}) \\
any & any & $D_f \cup D_g = \{p_1, \ldots, p_k\}$ & Coverage (Thm \ref{thm:CC}) \\
\hline
\end{tabular}
\end{center}

\begin{thebibliography}{99}

\bibitem{Birkhoff1940}
G.~Birkhoff,
\emph{Lattice Theory},
Amer.\ Math.\ Soc.\ Colloq.\ Publ.\ 25,
American Mathematical Society, 1940.

\bibitem{DrapalWanless2021}
A.~Dr\'apal and I.~M.~Wanless,
\emph{Maximally non-associative quasigroups},
J.\ Combin.\ Theory Ser.\ A \textbf{184} (2021), 105510.

\bibitem{DummitFoote}
D.~S.~Dummit and R.~M.~Foote,
\emph{Abstract Algebra}, 3rd ed.,
Wiley, 2004.

\bibitem{HardyWright}
G.~H.~Hardy and E.~M.~Wright,
\emph{An Introduction to the Theory of Numbers}, 6th ed.,
Oxford University Press, 2008.

\bibitem{IrelandRosen}
K.~Ireland and M.~Rosen,
\emph{A Classical Introduction to Modern Number Theory}, 2nd ed.,
Graduate Texts in Math.\ 84, Springer, 1990.

\bibitem{Lang2002}
S.~Lang,
\emph{Algebra}, 3rd ed.,
Graduate Texts in Math.\ 211, Springer, 2002.

\bibitem{Ore1942}
O.~Ore,
\emph{Theory of equivalence relations},
Duke Math.\ J.\ 9 (1942), 573--627.

\bibitem{Stanley2012}
R.~P.~Stanley,
\emph{Enumerative Combinatorics, Volume 1}, 2nd ed.,
Cambridge University Press, 2012.

\bibitem{SandersGish-CorrC}
B.~R.~Sanders and M.~Gish,
\emph{Corrected Theorem C: $n = 15$ Counterexample and the Sharpened
$M{+}A$ Condition},
submitted to \emph{J.\ Number Theory}, 2026 (companion).

\bibitem{SandersGish-Flatness}
B.~R.~Sanders and M.~Gish,
\emph{A Flatness Obstruction on Squarefree $\Z/n\Z$: Four Algebraic
Structures and the 4-Core Algebraic Center},
submitted to \emph{Algebraic Combinatorics}, 2026 (companion).

\end{thebibliography}

\end{document}
